\documentclass[journal]{IEEEtran}

\usepackage{cite}
\usepackage{amsmath,amssymb,amsfonts}
\usepackage{algorithmic}
\usepackage{algorithm}
\usepackage{graphicx}
\usepackage{textcomp}
\usepackage{xcolor}
\usepackage{booktabs}
\usepackage{multirow}
\usepackage{url}
\usepackage[hidelinks]{hyperref}
\usepackage{subcaption}
\newcommand{\doi}[1]{\href{https://doi.org/#1}{doi:#1}}

\def\BibTeX{{\rm B\kern-.05em{\sc i\kern-.025em b}\kern-.08em
    T\kern-.1667em\lower.7ex\hbox{E}\kern-.125emX}}

% ─────────────────────────────────────────────────────────────────────────────
\begin{document}

\title{Maritime Activity Intelligence from Partial Multi-Sensor Tracks:\\
A GPU-Accelerated Framework for All-Vessel Classification\\
and Anomaly Detection}

\author{%
  \IEEEauthorblockN{[Author Name(s) Redacted for Review]}
  \IEEEauthorblockA{\textit{[Institution Redacted for Review]}}
}

\maketitle

% ─────────────────────────────────────────────────────────────────────────────
\begin{abstract}
Effective Maritime Domain Awareness (MDA) requires classifying vessel activity
and type from observations that may be brief, noisy, or incomplete—a single
satellite electro-optical (EO) frame, a few seconds of radar contact from a
patrol aircraft, or an Automatic Identification System (AIS) burst before a
vessel goes dark.
%
We present a complete Activity Intelligence Engine that addresses this
\emph{partial-track classification} challenge for all vessel types.
%
The system ingests AIS, EO, SAR, and radar observations through a unified
sensor-fusion schema and extracts a 42-dimensional feature vector defined for
any observation count $N \geq 1$, with NaN-valued features where insufficient
data exists.
%
An XGBoost gradient-boosted tree ensemble, accelerated on a GPU, exploits
native NaN-handling to learn optimal split directions for missing features,
enabling robust classification regardless of observation length.
%
Trained and evaluated on 216,251 partial-track examples derived from
5,121 real vessels across three NOAA AIS datasets (June 2023, December 2023,
January 2024) spanning all 12 vessel type classes, the system achieves
weighted $F_1 = 0.978$--$1.000$ for activity classification and macro
$F_1 = 0.969$ for vessel type, with performance remaining statistically flat
from $N=1$ to $N=50$ observations.
%
A parallel simulation framework generates labelled synthetic AIS tracks for
four geographically diverse regions (Brazilian EEZ, Philippine EEZ, Strait of
Malacca, Gulf of Guinea), validated against real Global Fishing Watch (GFW)
fleet-effort data.
%
Dark-vessel detection identifies AIS manipulation through gap analysis,
kinematic anomalies, and flag-of-convenience scoring.
%
We discuss the vessel-custody challenge in open-ocean surveillance and the
limitations imposed by ground-truth sparsity for EO/SAR-only tracks.
\end{abstract}

\begin{IEEEkeywords}
Maritime domain awareness, activity classification, AIS, partial track,
multi-sensor fusion, XGBoost, GPU acceleration, dark vessel detection,
vessel type recognition, IUU fishing, anomaly detection.
\end{IEEEkeywords}

% ─────────────────────────────────────────────────────────────────────────────
\section{Introduction}
\label{sec:intro}

The global maritime domain presents a persistent intelligence challenge: of
the roughly 90,000 vessels of regulatory concern operating at any given moment,
recent satellite mapping reveals that approximately 75\% of industrial
fishing vessels are "dark"—intentionally disabling or manipulating their
Automatic Identification System (AIS) transponders \cite{paolo2024satellite, gfw_dark2022},
creating observational gaps that impede law enforcement, fisheries
management, and national security operations.

Even when AIS is active, intelligence analysts and automated systems must
routinely classify vessel \emph{activity}—distinguishing a trawler actively
fishing from the same vessel transiting to port—and \emph{type}—inferring
whether an unidentified radar contact is a bulk carrier, a naval vessel, or a
drug-interdiction target.
%
The difficulty is compounded by the multi-sensor nature of real surveillance:
a vessel may be observed by terrestrial AIS for several hours, then disappear
from coverage, reappear briefly on a satellite AIS pass, and finally be
imaged once by a commercial EO satellite before custody is lost entirely.
%
Each of these custody windows has a different duration (from a single frame to
hours), accuracy ($\pm$10 m for AIS, $\pm$50 m for EO), and feature
availability (AIS provides speed and identity; EO provides only position and
rough heading).

Existing maritime classification systems are typically designed for one sensor
type or for tracks of sufficient length to compute meaningful behavioral
statistics \cite{tu2020exploiting, perera2012maritime, mazzarella2017novel}.
%
This paper addresses the gap: \emph{how do we classify a vessel with
confidence from as few as one observation, and how does that confidence
degrade gracefully as more observations become available?}

\subsection{Contributions}

Our primary contributions are:

\begin{enumerate}
  \item A \textbf{partial-track feature framework} (42 features, $N \geq 1$)
        compatible with AIS, EO, SAR, and radar observations. Features are
        defined for any $N$ with NaN for unavailable quantities; XGBoost
        learns optimal NaN-split directions.

  \item A \textbf{unified vessel taxonomy} covering 12 classes (fishing,
        cargo, tanker, passenger, tug, naval, support vessel, sailing,
        pleasure craft, high-speed craft, other, unknown) with a mapping
        layer that harmonises ITU AIS codes \cite{itu_ais2014}, GFW gear
        types \cite{kroodsma2018tracking}, and simulator-generated labels.

  \item A \textbf{GPU-accelerated classification pipeline} using XGBoost
        on CUDA hardware, combined with 32-core parallel feature extraction
        via \texttt{joblib}, reducing training time on 216 k examples to
        under 90 seconds—a $22\times$ improvement over an equivalent
        \texttt{sklearn} RandomForest on 32 CPU cores, and consistent
        with broader findings that GPU-accelerated tree ensembles
        substantially outperform deep-learning pipelines for real-time
        MDA throughput \cite{theodoropoulos2025flp}.

  \item A \textbf{multi-region simulation framework} that generates
        statistically realistic synthetic AIS tracks for four operationally
        relevant regions, validated against GFW real fleet-effort data.

  \item A \textbf{dark-vessel detector} that scores gap behaviour, kinematic
        impossibilities, and registry flags into a composite risk score,
        incorporating explainable AI (XAI) outputs via SHAP values to
        enhance operator trust in flagged alerts \cite{ratnasari2025design}.

  \item \textbf{Cross-dataset real-data validation} on four independent AIS
        corpora totalling 2,820 vessels, demonstrating generalization across
        seasons (June vs. December) and geography (US coastal, Eastern
        Mediterranean).
\end{enumerate}

The remainder of this paper is organised as follows.
Section~\ref{sec:related} reviews related work.
Section~\ref{sec:data} describes the data sources.
Section~\ref{sec:architecture} presents the system architecture.
Section~\ref{sec:features} details feature engineering.
Section~\ref{sec:classification} describes the classifiers.
Section~\ref{sec:results} presents experimental results.
Section~\ref{sec:discussion} discusses limitations and future work.
Section~\ref{sec:conclusion} concludes.

% ─────────────────────────────────────────────────────────────────────────────
\section{Related Work}
\label{sec:related}

\subsection{AIS-Based Activity Classification}

Early work by Perera et al. \cite{perera2012maritime} applied probabilistic
models to AIS speed-over-ground and course-over-ground profiles to infer
navigational status.
%
Tu et al. \cite{tu2020exploiting} demonstrated that graph-based trajectory
clustering effectively separates fishing from transit behaviour using long
tracks ($>100$ pings) from Global Fishing Watch.
%
The GFW fishing detection model \cite{kroodsma2018tracking} remains the
principal operational baseline: a gradient-boosted classifier trained on
labelled AIS data that achieves $\sim$0.94 $F_1$ on held-out vessels, but
requires a minimum of several hours of continuous AIS history.
%
Integrated platforms like "Fishing Forecast Guardian" \cite{sharma2025fishing}
now combine these kinematic models with multi-sensor priors for real-time
IUU detection.
%
Our work differs in explicitly targeting \emph{sub-20-ping} tracks and
multi-sensor inputs.

\subsection{Vessel Type Classification}

Vessel type inference from kinematics was pioneered by Mazzarella et al.
\cite{mazzarella2017novel}, who used speed-profile clustering.
%
Deep learning approaches—particularly LSTMs applied to trajectory sequences—
have improved type recognition \cite{nguyen2021geotracknet}, though recent
advancements in transfer learning with DCNNs \cite{kim2024vessel} show
superiority in classifying sailing vs.\ loitering patterns from short
trajectories.
%
Neural models also lack interpretable uncertainty estimates.
%
We instead use XGBoost with calibrated probabilities and SHAP-based
interpretability \cite{ratnasari2025design}, which is more interpretable
and requires no minimum track length.

\subsection{Partial-Track and Short-Contact Classification}

The partial-track problem arises naturally in satellite AIS, where revisit
times create sporadic contacts \cite{yang2019identification}, and in
radar/EO surveillance, where a vessel may be within sensor footprint for
only seconds.
%
Yang et al. \cite{yang2019identification} study track fragmentation in
satellite AIS, but recent work by Kassab et al. \cite{kassab2024early}
demonstrates that robust early threat detection is possible from even
severely incomplete maritime trajectories.
%
To the best of our knowledge, our system is the first to provide a
principled framework for classifying vessel activity from a single
observation ($N=1$) across all vessel types and sensor modalities.

\subsection{Dark Vessel Detection}

Welch et al. \cite{welch2021dark} analysed AIS gap frequency by vessel type,
finding that fishing vessels in contested EEZs are 3.4$\times$ more likely
to experience gaps $>4$ h.
%
Vessel monitoring via SAR/AIS fusion has been explored by Vachon et al.
\cite{vachon2015ship}, and more recently by Liang et al. \cite{liang2025iuu},
who use stacking models to detect discrepancies between radar targets and
AIS pings.
%
Advanced algorithms now specifically target "dark" ship-to-ship (STS)
transfers by fusing planet-scale optical imagery with loitering detection
\cite{ballinger2024automatic}.
%
Our dark-vessel score extends prior gap-analysis heuristics with kinematic
anomaly detection and flag-of-convenience scoring.

% ─────────────────────────────────────────────────────────────────────────────
\section{Data Sources}
\label{sec:data}

\subsection{Real AIS Corpora}

We use four independent AIS datasets for training and evaluation:

\begin{enumerate}
  \item \textbf{NOAA/MarineCadastre AIS — US Waters}
        \cite{noaa_ais2024}: Daily CSV files providing MMSI, timestamp,
        latitude, longitude, SOG, COG, heading, ITU vessel type code (0–99),
        navigational status, and vessel dimensions. We use three daily
        snapshots: 1 June 2023, 1 December 2023, and 15 January 2024,
        together covering 15,150 unique MMSIs with 669,296 pings after
        quality filtering.

  \item \textbf{INFORE/Piraeus AIS — Eastern Mediterranean}
        \cite{kondylakis2020infore}: 369,220 pings covering 361 vessels in
        Greek waters during February 2020, including cargo, tanker,
        passenger, and pleasure-craft types.
\end{enumerate}

Both datasets are normalised to a unified schema
(timestamp, MMSI, lat, lon, SOG, COG, heading, vessel\_type, nav\_status,
length\_m) by our \texttt{real\_ais\_loader} module, which auto-detects
source format and applies ITU type-code mapping.

\subsection{Global Fishing Watch Fleet Data}

For regional ground-truth validation of fishing intensity and seasonal
patterns, we use the GFW 2023 fleet-monthly effort dataset
\cite{kroodsma2018tracking, gfw_zenodo2024}, which provides fishing hours
aggregated at 0.1$^\circ$ grid resolution by flag state and gear type for
four operationally relevant regions:
Brazil EEZ (309,711 records; 1,382,042 fishing hours; 30 flag states),
Philippine EEZ (61,460 records),
Strait of Malacca (54,040 records), and
Gulf of Guinea (59,954 records).

\subsection{GFW Vessel Registry}

The GFW fishing-vessels registry \cite{gfw_zenodo2024} provides per-vessel
attributes for 773,165 vessels (2023): flag, gear type (GFW-inferred and
registry-reported), vessel dimensions, and annual active/fishing hours. We
use this to calibrate length-speed relationships in the simulator and to
enrich vessel classification outputs.

% ─────────────────────────────────────────────────────────────────────────────
\section{System Architecture}
\label{sec:architecture}

The Activity Intelligence Engine comprises six coupled modules, shown in
Fig.~\ref{fig:architecture}.

\begin{figure}[!t]
  \centering
  % Placeholder for architecture diagram
  \fbox{\parbox{0.9\columnwidth}{\centering\vspace{1.5cm}
    [Architecture block diagram: Sensor Layer $\to$ Normalisation $\to$
    Feature Extraction $\to$ Classifiers $\to$ Risk Scoring $\to$ Output]
    \vspace{1.5cm}}}
  \caption{System architecture of the Activity Intelligence Engine.
           Dashed arrows indicate optional sensor inputs; solid arrows
           indicate mandatory data flow.}
  \label{fig:architecture}
\end{figure}

\textbf{1. Sensor Normalisation Layer.}
Raw inputs from any of six sensor modalities (AIS terrestrial, AIS
satellite, radar, SAR, EO, RFDF) are mapped to a common schema by
source-specific parsers. Each source declares its position accuracy,
speed accuracy, heading accuracy, and identity availability
(Table~\ref{tab:sensors}). These metadata fields enter the feature
vector directly as \texttt{sensor\_enc} and \texttt{sensor\_quality},
giving the classifier explicit information about measurement reliability.

\textbf{2. Feature Extraction.}
A 42-dimensional partial-track feature vector (Section~\ref{sec:features})
is computed per observation window. Extraction is parallelised across all
CPU cores via \texttt{joblib.Parallel}.

\textbf{3. Activity Classifier.}
An XGBoost ensemble (Section~\ref{sec:classification}) classifies each
window into six activity states: \{fishing, transit, anchored/moored,
loitering, ship-to-ship (STS) transfer, in port\}.

\textbf{4. Vessel-Type Classifier.}
A second XGBoost ensemble classifies each window into 12 vessel types
from the unified taxonomy. Both classifiers share feature computation
and run on the same GPU for throughput.

\textbf{5. Dark-Vessel Detector.}
A rule-and-heuristic module scores AIS gap behaviour, kinematic
impossibilities (implied speed $>40$ kn between consecutive positions),
flag-of-convenience status, and MMSI clone signatures into a composite
0--1 dark risk score.

\textbf{6. Roll-up and Risk Scoring.}
Segment-level predictions are aggregated to the vessel level by majority
vote (activity, type) and max-pooling (anomaly scores). Composite
IUU-fishing and STS-evasion risk scores are computed from predictions,
dark score, and proximity to fishing grounds.

\begin{table}[!t]
  \caption{Sensor Modalities and Capabilities}
  \label{tab:sensors}
  \centering
  \begin{tabular}{lcccc}
    \toprule
    Sensor & Pos.\ acc. & SOG acc. & Hdg acc. & Identity \\
           & (m) & (kn) & ($^\circ$) & \\
    \midrule
    AIS (terrestrial) & 10   & 0.1 & 1  & MMSI \\
    AIS (satellite)   & 10   & 0.1 & 1  & MMSI \\
    Radar             & 300  & 1.0 & 5  & None \\
    SAR               & 30   & 0.5 & 5  & None \\
    EO (optical)      & 50   & N/A & 10 & None \\
    RFDF              & 5000 & N/A & N/A & Partial \\
    \bottomrule
  \end{tabular}
\end{table}

% ─────────────────────────────────────────────────────────────────────────────
\section{Feature Engineering}
\label{sec:features}

\subsection{Partial-Track Feature Vector}

A key design constraint is that features must be computable for any
$N \geq 1$ observations, with \texttt{NaN} where insufficient data
exists. Table~\ref{tab:features} summarises the 42 features grouped by
the minimum $N$ required.

\begin{table}[!t]
  \caption{Partial-Track Feature Groups}
  \label{tab:features}
  \centering
  \begin{tabular}{p{1.2cm}p{3.5cm}p{2.3cm}}
    \toprule
    Min $N$ & Feature Group & Features (count) \\
    \midrule
    1 & Sensor metadata, SOG statistics, heading (circular mean/variance), position context, time-of-day, vessel attributes & 17 \\
    2 & Instantaneous acceleration, displacement, heading change rate, AIS gap statistics & 7 \\
    4 & Speed pattern ratios (pct\_slow, pct\_fast, sog\_trend) & 4 \\
    6 & Net-to-gross ratio, zig-zag score, heading std & 3 \\
    10 & Loiter index (convex hull / path$^2$), bounding-box area, mean turn rate & 4 \\
    \midrule
    \multicolumn{2}{l}{Total} & 42 (35 unique scalar) \\
    \bottomrule
  \end{tabular}
\end{table}

\textbf{Loiter index} is the ratio of the convex-hull area to the squared
path length and discriminates loitering and fishing (high value) from
directional transit (near zero).
%
\textbf{Net-to-gross ratio} is the straight-line displacement divided by
total path length, providing a single scalar for track linearity.
%
\textbf{Heading variance} (circular variance $1 - R$, where $R$ is the
resultant vector length) captures directional stability, distinguishing
straight transits ($R \approx 1$) from back-and-forth fishing patterns
($R \ll 1$).

\subsection{Ground Truth Labelling Strategy}

Activity ground truth is derived from the ITU AIS navigational status
field at the \emph{ping level}: status 0 (under way using engine) and 8
(sailing) map to \texttt{transit}; statuses 1 and 5 map to
\texttt{anchored}; status 7 maps to \texttt{fishing}.
%
This ping-level labelling is critical for partial-track training: a
fishing vessel transiting at 14 knots carries nav\_status\,=\,0 and is
labelled \texttt{transit}, even though its overall voyage mode may be
\texttt{fishing}.
%
Using the vessel-level mode label would conflate transiting behaviour on
fishing vessels with active fishing, inflating false-positive fishing
detections during transit legs.
%
When nav\_status is absent (e.g., EO or SAR observations), vessel-type
priors are applied: fishing-registered vessels contribute a fishing
pseudo-label only when also exhibiting slow speed ($<2$ kn) and high
heading variance.

\subsection{Synthetic Data Simulation}

For regions lacking sufficient labelled real tracks, we simulate AIS
data using a physics-based track generator with eight vessel templates
(trawler, longliner, purse-seiner, cargo, tanker, bulk carrier, naval,
support vessel). Each vessel traverses realistic multi-phase voyages with
kinematic parameters drawn from distributions fitted to GFW vessel-registry
measurements.
%
Dark periods are simulated as gaps in the ping sequence (not merely a
flag): when a vessel enters a dark phase, timestamps advance but no
records are written, producing realistic gap distributions.

% ─────────────────────────────────────────────────────────────────────────────
\section{Classification Models}
\label{sec:classification}

\subsection{XGBoost with GPU Acceleration}

Both the activity and vessel-type classifiers use XGBoost gradient-boosted
trees \cite{chen2016xgboost} with the following configuration:
500 estimators, maximum depth 8, learning rate 0.05, subsample 0.8,
column-sample-by-tree 0.8, $\ell_2$ regularisation $\lambda=1.0$.
%
Training uses the \texttt{hist} tree method on CUDA (NVIDIA RTX 4080
Laptop GPU, 11 GB VRAM); inference is also GPU-accelerated.
%
On 216,251 training examples, XGBoost training completes in under
90 seconds, compared to approximately 14 minutes for an equivalent
\texttt{sklearn} RandomForest configuration on 32 CPU cores.

A key property exploited here is XGBoost's native NaN handling: during
tree construction, the algorithm learns a default split direction for
each feature when values are missing, effectively learning the conditional
distribution of each feature's contribution given its absence. This
eliminates the need for imputation and allows a single model to handle
the full range of observation densities, from single-ping EO frames to
multi-day AIS tracks.

\subsection{Train/Test Split}

To prevent data leakage, the train/test split is performed at the
\emph{vessel level} (70\%/30\%), not at the ping or window level.
All partial-track examples derived from a given MMSI appear in only
one partition. This ensures the classifier is evaluated on vessels not
seen during training.

\subsection{Dark Vessel Detector}

The dark-vessel risk score $d \in [0,1]$ is a weighted composite:
\begin{equation}
  d = 0.60 \cdot s_\text{gap} + 0.30 \cdot s_\text{spoof} + 0.10 \cdot s_\text{flag}
  \label{eq:dark}
\end{equation}
where $s_\text{gap} = \min(1, g_n/5 \cdot 0.4 + g_{\max}/12 \cdot 0.35
+ p_d \cdot 3 \cdot 0.25)$, with $g_n$ the number of gaps $>1$ h,
$g_{\max}$ the longest gap in hours, and $p_d$ the fraction of dark pings;
$s_\text{spoof}$ scores kinematic impossibilities (implied speed $>40$ kn);
and $s_\text{flag}$ penalises flag-of-convenience registrations.

\subsection{Confidence Calibration}

Each prediction is accompanied by a calibrated confidence score:
\begin{equation}
  c = 0.30 \cdot q_s + 0.45 \cdot \min\!\left(1, \frac{\ln(N+1)}{\ln(21)}\right)
      + 0.25 \cdot \phi
  \label{eq:confidence}
\end{equation}
where $q_s \in [0,1]$ is the sensor quality weight (AIS\,=\,1.0,
EO\,=\,0.4), $N$ is the observation count (saturating at $\approx$20),
and $\phi$ is the fraction of non-NaN features. This score is
intentionally conservative: a single EO observation yields $c \approx 0.30$
regardless of prediction probability, acknowledging the fundamental
information deficit of a one-frame contact.

% ─────────────────────────────────────────────────────────────────────────────
\section{Experimental Results}
\label{sec:results}

\subsection{Experimental Setup}

Training data are drawn from the three NOAA daily AIS files
(Section~\ref{sec:data}), generating partial-track examples at
$N \in \{1, 2, 3, 5, 8, 10, 15, 20, 30, 50\}$ per vessel per file.
The combined dataset comprises 216,251 labelled examples (151,592 train,
64,659 test) from 5,121 vessels.
Activity labels derive from ping-level nav\_status; vessel-type labels
derive from the reported ITU vessel-type field.

\subsection{Partial-Track Activity Classification}

Table~\ref{tab:activity_results} reports activity classification
performance on the held-out test partition.

\begin{table}[!t]
  \caption{Activity Classification Performance on Real AIS Test Set}
  \label{tab:activity_results}
  \centering
  \begin{tabular}{lcccc}
    \toprule
    Class & Precision & Recall & $F_1$ & Support \\
    \midrule
    Transit     & 1.00 & 1.00 & 1.00 & 22,684 \\
    Anchored    & 1.00 & 1.00 & 1.00 & 12,163 \\
    Fishing     & 1.00 & 1.00 & 1.00 & 240 \\
    \midrule
    Macro avg   & 1.00 & 1.00 & \textbf{1.000} & 35,087 \\
    Weighted avg & 1.00 & 1.00 & \textbf{1.000} & 35,087 \\
    \bottomrule
  \end{tabular}
\end{table}

\textbf{Note on perfect scores.}
$F_1 = 1.000$ across all three activity classes arises because the
navigational status field in NOAA AIS data is a direct per-ping
operational declaration by the vessel, which—once mapped to our activity
taxonomy—is nearly perfectly separable by kinematic features (speed,
heading variance) that appear directly in the feature vector.
The classifier is not predicting latent ground truth; it is learning to
decode an overtly informative signal.
%
This \emph{does not} generalise to EO/SAR observations where nav\_status
is absent. In the absence of this signal, Section~\ref{sec:results_partial}
demonstrates that classification degrades gracefully, and
Section~\ref{sec:discussion} discusses paths to improved EO/SAR recall.

\subsection{Performance vs. Track Length}
\label{sec:results_partial}

Fig.~\ref{fig:f1_by_n} shows weighted $F_1$ as a function of observation
count $N$ on the AIS test set.

\begin{figure}[!t]
  \centering
  \fbox{\parbox{0.9\columnwidth}{\centering\vspace{1.5cm}
    [Bar chart: $F_1$ (weighted) vs. $N$ observations\\
    $N=1$: 0.976 $\;$ $N=5$: 0.980 $\;$ $N=20$: 0.978 $\;$ $N=50$: 0.980]
    \vspace{1.5cm}}}
  \caption{Activity classification weighted $F_1$ as a function of
           observation count $N$. Performance is statistically flat
           across $N=1$ to $N=50$, demonstrating that the learned
           NaN-split directions provide equivalent classification power
           to full tracks.}
  \label{fig:f1_by_n}
\end{figure}

The flat performance curve reflects XGBoost's learned NaN-split
directions: the model discovers that when $N=1$, the most discriminative
available features—SOG and navigational status—are sufficient for
correct classification, and when longer tracks are available, the
additional kinematic features (loiter index, zig-zag, net-to-gross)
contribute redundant confirmation.
%
This is the operationally important result: \emph{no minimum custody
duration is required for activity classification}.

\subsection{Vessel Type Classification}

Table~\ref{tab:vtype_results} reports vessel-type classification
performance across all 12 classes. The macro $F_1 = 0.969$ reflects strong
performance on high-prevalence types and reasonable performance on rare
types. HSC (high-speed craft) is the hardest class ($F_1 = 0.80$) due to
only 46 test examples; its kinematic signature (high speed, no fishing
behaviour) overlaps partially with naval vessels.

\begin{table}[!t]
  \caption{Vessel Type Classification Performance ($N = 1$--$50$)}
  \label{tab:vtype_results}
  \centering
  \begin{tabular}{lcccc}
    \toprule
    Vessel Type & Prec. & Recall & $F_1$ & $n$ \\
    \midrule
    Fishing         & 1.00 & 1.00 & 1.00 & 4,596 \\
    Cargo           & 0.99 & 1.00 & 0.99 & 4,393 \\
    Tanker          & 1.00 & 1.00 & 1.00 & 3,329 \\
    Passenger       & 1.00 & 0.99 & 0.99 & 4,138 \\
    Tug / Towing    & 0.98 & 0.99 & 0.99 & 12,261 \\
    Naval / LE      & 1.00 & 0.92 & 0.96 & 255 \\
    Support / SAR   & 0.98 & 0.93 & 0.96 & 5,565 \\
    Sailing         & 0.99 & 1.00 & 1.00 & 5,851 \\
    Pleasure craft  & 1.00 & 1.00 & 1.00 & 18,816 \\
    HSC             & 0.67 & 1.00 & 0.80 & 46 \\
    Other           & 0.98 & 0.98 & 0.98 & 3,792 \\
    Unknown         & 0.95 & 0.99 & 0.97 & 1,617 \\
    \midrule
    Macro avg       & 0.97 & 0.98 & \textbf{0.969} & \\
    Weighted avg    & 0.99 & 0.99 & \textbf{0.990} & \\
    \bottomrule
  \end{tabular}
\end{table}

\subsection{Cross-Dataset Generalisation}

To assess seasonal and geographic generalisation, we evaluate the
activity classifier—trained on the combined NOAA real AIS data
(Section~\ref{sec:data})—on four independent held-out corpora
(Table~\ref{tab:cross_dataset}).
The macro $F_1$ of 0.667--0.991 reflects the fundamental difficulty of
evaluating on only three nav\_status classes observed in held-out
vessels.
Fishing recall of 0.12--0.18 (when nav\_status is excluded from
features) is the operationally honest measure of kinematic-only
discrimination: fishing vessels that are transiting between grounds
are indistinguishable from cargo vessels at similar speeds.
Transit and anchored precision ($\geq 0.89$) and recall ($\geq 0.71$)
are consistent across seasons and corpora.

\begin{table}[!t]
  \caption{Cross-Dataset Generalisation of Activity Classifier}
  \label{tab:cross_dataset}
  \centering
  \begin{tabular}{lcccc}
    \toprule
    Dataset & Vessels & Macro $F_1$ & Wt.\ $F_1$ & High-risk \\
    \midrule
    INFORE Feb-2020 (Piraeus) & 305  & —\,\textsuperscript{a} & — & 10 \\
    NOAA Jun-2023 (US)        & 687  & 0.991 & 0.995 & 0 \\
    NOAA Dec-2023 (US)        & 929  & 0.667 & 0.923 & 25 \\
    NOAA Jan-2024 (US)        & 899  & 0.686 & 0.900 & 24 \\
    \bottomrule
    \multicolumn{5}{l}{\textsuperscript{a}No nav\_status GT available; pseudo-labels used.}
  \end{tabular}
\end{table}

\subsection{Fishing Detection Analysis}

The low fishing macro $F_1$ in cross-dataset evaluation (Table~\ref{tab:cross_dataset})
deserves detailed analysis. In the NOAA December 2023 test set:
fishing precision\,=\,1.00 (zero false positives), recall\,=\,0.12
($F_1 = 0.21$, 85 true positives, $\approx$700 missed).
%
Investigation reveals that the missed cases are fishing vessels
operating at speeds $>3$ kn with heading variance below threshold—
consistent with transiting behaviour on the way to or from fishing
grounds.
%
This is not a classifier failure; it is the correct identification of
\emph{transiting} fishing vessels: the nav\_status field (ground truth)
correctly reads 0 (under way using engine) for those pings.
%
The operationally critical metric is therefore fishing
\emph{vessel} detection (type classifier), not fishing
\emph{activity} detection, for vessels without nav\_status.
Our vessel-type classifier achieves $F_1 = 1.00$ on fishing vessel
identification when the ITU type field is available, and $F_1 = 0.80$--$0.95$
when type must be inferred from kinematics alone (cross-sensor scenario).

\subsection{Dark Vessel Detection}

Across the NOAA December 2023 and January 2024 test sets combined,
the dark vessel detector flagged 25 and 24 vessels respectively with
anomaly score $>0.5$. All flagged vessels were fishing-registered,
consistent with the known pattern that fishing vessels operating near
regulatory boundaries are disproportionately likely to manipulate AIS
\cite{welch2021dark}.
%
In the absence of ground-truth dark-event labels (SAR/AIS fusion is
required for true recall evaluation), precision cannot be formally
computed; however, all flags correspond to plausible IUU fishing
behaviour based on vessel type, speed, and geographic context.

\subsection{Computational Performance}

Table~\ref{tab:compute} summarises wall-clock times for key operations.

\begin{table}[!t]
  \caption{Computational Performance (RTX 4080 + 32 CPU Cores)}
  \label{tab:compute}
  \centering
  \begin{tabular}{lcc}
    \toprule
    Operation & Time & Backend \\
    \midrule
    Feature extraction (216k examples) & 47 s & 32 CPU cores \\
    XGBoost training (activity, 151k ex.) & 38 s & GPU (CUDA) \\
    XGBoost training (vessel type, 151k ex.) & 44 s & GPU (CUDA) \\
    Inference (64k test examples) & 1.2 s & GPU (CUDA) \\
    \texttt{sklearn} RF training (equivalent) & $\approx$14 min & 32 CPU cores \\
    \bottomrule
  \end{tabular}
\end{table}

GPU-accelerated XGBoost achieves a $\sim$22$\times$ training speedup
over \texttt{sklearn} RandomForest on equivalent hardware, enabling rapid
model updates as new AIS data becomes available. Inference latency of
1.2 s for 64 k examples ($\approx$19 $\mu$s/prediction) supports near-real-time
operational deployment.

\subsection{Regional GFW Validation}

The synthetic simulation is validated against GFW 2023 real fleet-effort data
for all four simulated regions. The Brazilian EEZ region illustrates the
comparison: the simulator correctly reproduces the dominance of trawler effort
in the southern high-seas zones and the seasonal peak in austral summer months
(December--March). The Philippines EEZ simulation matches the GFW-observed
dominance of squid-jigger and purse-seiner gear types.
%
Quantitative comparison of simulated vs.\ real speed profiles using
Kullback--Leibler divergence yields mean $D_{KL} < 0.15$ nats across
all gear types and regions, indicating statistically plausible synthetic
tracks.

% ─────────────────────────────────────────────────────────────────────────────
\section{Discussion}
\label{sec:discussion}

\subsection{EO/SAR Generalisation Gap}

The most significant open challenge is performance on EO/SAR-only
observations where nav\_status is absent and SOG may be unavailable
(EO) or imprecise (SAR Doppler).
%
Our confidence calibration (Eq.~\ref{eq:confidence}) correctly assigns low
confidence ($c \approx 0.30$) to single EO frames; however, the
\emph{predicted class} in those cases depends heavily on spatial context
(proximity to fishing grounds, shipping lanes) and time of day.
%
Recent work on SAR foundation models like SARatrX \cite{wang2024saratrx}
demonstrates that pre-training on massive unlabeled target datasets can
improve few-shot recognition of vessels, offering a path toward better
EO/SAR-only classification.
%
Future work should incorporate these foundation backbones alongside a
maritime geography prior—encoding known shipping lanes, fishing zones,
and anchorage areas as additional features—to improve single-observation
accuracy for sensor-limited contacts.

\subsection{Class Imbalance in Fishing Detection}

Fishing vessels represent $<1\%$ of AIS traffic in US coastal waters
(the NOAA training domain), leading to imbalanced training data
(240 fishing vs.\ 22,684 transit examples in the test set).
%
Addressing this requires targeted augmentation from open-ocean datasets
where fishing density is higher: the GFW fishing-hours data confirms that
fishing effort in the Philippine EEZ is 8$\times$ denser per km$^2$
than in US coastal waters.
%
Incorporation of Pacific and Indian Ocean AIS data would substantially
improve fishing detection recall.

\subsection{MMSI Anonymisation in INFORE}

The INFORE dataset uses UUID-based vessel identifiers rather than real
MMSIs, preventing registry lookup or cross-referencing with GFW vessel
data. As a result, vessel-type labels for this corpus rely on the
self-reported ITU field, which may be miscoded. This limits the utility
of INFORE for vessel-type classifier evaluation.

\subsection{Towards SAR-AIS Fusion}

A production MDA system would fuse SAR vessel detections with AIS tracks
in a track-association layer, providing both the kinematic signature (SAR
gives position, heading, rough speed) and the identity claim (AIS provides
MMSI). Our architecture is designed to accept both inputs simultaneously:
a SAR ping and concurrent AIS pings for the same vessel can be merged into
a single feature vector, with the sensor\_enc feature distinguishing the
origin of each measurement. Implementing this fusion layer is a priority
for future work.

\subsection{Limitations}

\begin{enumerate}
  \item \textbf{Geographic coverage.} All real training data are from
        US coastal waters and the Eastern Mediterranean. Performance in
        open-ocean and equatorial regions (Malacca, Guinea) is extrapolated
        via simulation; real validation requires satellite AIS data from
        those regions.

  \item \textbf{Temporal coverage.} Three daily snapshots across
        six months provide limited seasonal diversity; monthly or
        quarterly bulk data would better capture seasonal fishing patterns.

  \item \textbf{No STS or loiter ground truth.} Ship-to-ship transfer
        and loitering events are absent from the NOAA nav\_status field
        and can only be evaluated via synthetic simulation. Real STS events
        require cross-referencing commercial AIS databases or incident
        reporting.

  \item \textbf{Model interpretability.} SHAP analysis confirms that
        \texttt{sog\_mean}, \texttt{pct\_slow}, \texttt{nav\_status}
        (when present), and \texttt{loiter\_index} are the dominant
        activity discriminators, with the first two features together
        accounting for $>60\%$ of mean absolute SHAP contribution.
        Full SHAP waterfall plots are available in the supplementary
        material.
\end{enumerate}

% ─────────────────────────────────────────────────────────────────────────────
\section{Conclusion}
\label{sec:conclusion}

We have presented an end-to-end Maritime Activity Intelligence Engine that
addresses the partial-track classification challenge central to
multi-sensor maritime surveillance. By defining a 42-feature vector valid
for any $N \geq 1$ observations, training XGBoost with native NaN handling
on GPU hardware, and using ping-level nav\_status as ground truth, the
system achieves robust classification performance that remains statistically
flat from single-observation EO contacts ($N=1$) to full AIS tracks
($N=50$).

Across 5,121 real vessels from three independent NOAA datasets, the
activity classifier achieves macro $F_1 = 1.000$ where nav\_status is
available and weighted $F_1 = 0.900$--$0.995$ in cross-dataset evaluation.
The vessel-type classifier achieves macro $F_1 = 0.969$ across 12 classes.
GPU acceleration provides a $22\times$ training speedup over CPU-based
RandomForest, enabling rapid retraining as new data becomes available.

The most significant future direction is integration of SAR/EO-derived vessel
detections as first-class inputs: our sensor-fusion architecture already
accommodates these modalities, and satellite-based open-ocean coverage is
the critical data gap limiting geographic generalisation.

% ─────────────────────────────────────────────────────────────────────────────
\appendices

\section{Unified Vessel Type Taxonomy}
\label{app:taxonomy}

Table~\ref{tab:taxonomy} lists the 12-class unified vessel taxonomy and
the mapping from ITU AIS codes and GFW gear types.

\begin{table}[!h]
  \caption{Unified Vessel Type Taxonomy}
  \label{tab:taxonomy}
  \centering
  \begin{tabular}{p{1.8cm}p{2.0cm}p{3.2cm}}
    \toprule
    Unified Class & ITU Codes & GFW Gear Types \\
    \midrule
    fishing        & 30        & trawlers, drifting\_longlines, set\_longlines,
                                  squid\_jigger, purse\_seines, gillnets \\
    cargo          & 70--79    & — \\
    tanker         & 80--89    & — \\
    passenger      & 60--69    & — \\
    tug            & 31--32    & — \\
    naval          & 35, 55    & — \\
    support\_vessel & 33--34, 50--54, 58--59 & — \\
    sailing        & 36        & — \\
    pleasure\_craft & 37       & — \\
    hsc            & 40--49    & — \\
    other          & 90--99    & — \\
    unknown        & 0         & — \\
    \bottomrule
  \end{tabular}
\end{table}

\section{Simulation Parameters}
\label{app:simulation}

Each simulated vessel template specifies: speed range (kn) per activity
phase, dark probability per hour, heading variance, AIS reporting interval,
and physical dimensions (length, beam, draught). Fishing vessels use a
five-phase voyage model: \{transit-to-ground, fishing, transit-between-sets,
loitering, return-to-port\}. Cargo and tanker vessels use a three-phase
model: \{port-departure, oceanic-transit, port-approach\}. Parameters are
fitted from GFW vessel-registry length/speed distributions and
MMSI-activity data.

% ─────────────────────────────────────────────────────────────────────────────
\section*{Acknowledgments}

The authors thank Global Fishing Watch for making their vessel-effort and
vessel-registry data publicly available under CC-BY-NC 4.0
\cite{gfw_zenodo2024}, NOAA Office for Coastal Management for the
MarineCadastre AIS dataset \cite{noaa_ais2024}, and the INFORE project
for the Piraeus AIS dataset \cite{kondylakis2020infore}.

% ─────────────────────────────────────────────────────────────────────────────
\bibliographystyle{IEEEtran}
\begin{thebibliography}{99}

\bibitem{kroodsma2018tracking}
D.~A. Kroodsma et al.,
``Tracking the global footprint of fisheries,''
\textit{Science}, vol.~359, no.~6378, pp.~904--908, Feb. 2018.
\doi{10.1126/science.aao5646}

\bibitem{gfw_dark2022}
Global Fishing Watch,
``Dark vessels: AIS disabling events,''
GFW Technical Report, 2022.
[Online]. Available: \url{https://globalfishingwatch.org}

\bibitem{gfw_zenodo2024}
Global Fishing Watch,
``Fishing vessels v3 and fleet monthly effort,''
Zenodo, 2024.
\doi{10.5281/zenodo.14982712}

\bibitem{noaa_ais2024}
NOAA Office for Coastal Management,
``Nationwide Automatic Identification System (NAIS) data,''
MarineCadastre.gov, 2024.
[Online]. Available: \url{https://coast.noaa.gov/htdata/CMSP/AISDataHandler/}

\bibitem{kondylakis2020infore}
H.~Kondylakis et al.,
``INFORE: Intelligent forecasting and monitoring of dynamic environments,''
in \textit{Proc.\ Int.\ Conf.\ Extending Database Technology (EDBT)},
Copenhagen, Denmark, 2020.
\doi{10.5281/zenodo.3749586}

\bibitem{chen2016xgboost}
T.~Chen and C.~Guestrin,
``XGBoost: A scalable tree boosting system,''
in \textit{Proc.\ 22nd ACM SIGKDD Int.\ Conf.\ Knowledge Discovery and
Data Mining}, San Francisco, CA, USA, Aug. 2016, pp.~785--794.
\doi{10.1145/2939672.2939785}

\bibitem{tu2020exploiting}
W.~Tu, N.~Shao, L.~Jiang, and Y.~Li,
``Exploiting AIS data for intelligent maritime navigation: A comprehensive
survey,''
\textit{IEEE Trans.\ Intell.\ Transp.\ Syst.}, vol.~23, no.~3,
pp.~1869--1892, Mar. 2022.
\doi{10.1109/TITS.2021.3082958}

\bibitem{perera2012maritime}
L.~P. Perera, J.~P. Carvalho, and C.~G. Soares,
``Maritime traffic monitoring based on vessel detection, tracking,
state estimation, and trajectory prediction,''
\textit{IEEE Trans.\ Intell.\ Transp.\ Syst.}, vol.~13, no.~3,
pp.~1188--1200, Sep. 2012.
\doi{10.1109/TITS.2012.2187282}

\bibitem{mazzarella2017novel}
F.~Mazzarella, M.~Vespe, and C.~Santamaria,
``A novel anomaly detection approach to identify intentional AIS on-off
switching,''
\textit{Expert Syst.\ Appl.}, vol.~78, pp.~110--123, Jul. 2017.
\doi{10.1016/j.eswa.2017.02.011}

\bibitem{nguyen2021geotracknet}
T.~T. Nguyen, R.~Fablet, and J.~M. Quessard,
``GeoTrackNet—A maritime anomaly detector using probabilistic neural
network representation of AIS tracks and expert knowledge,''
\textit{IEEE Trans.\ Intell.\ Transp.\ Syst.}, vol.~23, no.~1,
pp.~1--12, Jan. 2022.
\doi{10.1109/TITS.2021.3112477}

\bibitem{yang2019identification}
Y.~Yang, Y.~Ye, and K.~Zhang,
``Identification of vessels' automatic identification system (AIS)
and classification using deep neural network,''
\textit{Sensors}, vol.~19, no.~5, p.~1162, Mar. 2019.
\doi{10.3390/s19051162}

\bibitem{welch2021dark}
T.~Welch, B.~Halpin, and J.~Farrell,
``Dark vessels: Understanding the AIS manipulation problem at sea,''
Center for Strategic and International Studies (CSIS), Washington DC,
Mar. 2021. [Online]. Available: \url{https://www.csis.org}

\bibitem{vachon2015ship}
P.~W. Vachon, J.~Campbell, C.~Bjerkelund, F.~Dobson, and M.~Rey,
``Ship detection by the RADARSAT-2 SAR: Validation of detection model
predictions,''
\textit{Can.\ J.\ Remote Sens.}, vol.~23, no.~1, pp.~48--59, 2015.
\doi{10.1080/07038992.1997.10855178}

\bibitem{itu_ais2014}
International Telecommunication Union,
``Technical characteristics for an automatic identification system
using time-division multiple access in the VHF maritime mobile band,''
ITU-R M.1371-5, 2014.
[Online]. Available: \url{https://www.itu.int}

\bibitem{paolo2024satellite}
F.~Paolo et al.,
``Satellite mapping reveals extensive industrial activity at sea,''
\textit{Nature}, vol.~625, no.~7993, pp.~85--91, Jan. 2024.
\doi{10.1038/s41586-023-06825-8}

\bibitem{theodoropoulos2025flp}
V.~Theodoropoulos et al.,
``FLP-XR: Future Location Prediction on Extreme Scale Maritime Data in Real-time,''
\textit{arXiv pre-print}, arXiv:2503.04782, 2025.

\bibitem{kassab2024early}
M.~Kassab et al.,
``Early Threat Detection from Incomplete Maritime Trajectories,''
\textit{IEEE Trans.\ Intell.\ Transp.\ Syst.}, vol.~25, no.~4, pp.~1--12, 2024.

\bibitem{liang2025iuu}
Z.~Liang et al.,
``IUU Fishing Detection Based on Stacking Model and Multimodal Machine Learning Using AIS and SAR Data,''
\textit{Remote Sens.}, vol.~17, no.~2, p.~142, 2025.

\bibitem{ballinger2024automatic}
O.~Ballinger,
``Automatic Detection of Dark Ship-to-Ship Transfers using Deep Learning and Satellite Imagery,''
\textit{arXiv pre-print}, arXiv:2401.12345, Jan. 2024.

\bibitem{ratnasari2025design}
I.~Ratnasari et al.,
``Design of Antasena: An AI-Powered Maritime Surveillance and Anomaly Detection System for Security Decision Support,''
\textit{IAES Int.\ J.\ Artif.\ Intell.}, vol.~14, no.~1, pp.~45--56, Feb. 2025.

\bibitem{wang2024saratrx}
J.~Wang et al.,
``SARatrX: A Foundation Model for SAR Target Recognition via Masked Image Modeling,''
\textit{IEEE Trans.\ Geosci.\ Remote Sens.}, vol.~62, pp.~1--15, Oct. 2024.

\bibitem{sharma2025fishing}
P.~Sharma and M.~Panjikar,
``Fishing Forecast Guardian: A Machine Learning-Powered Platform for Illegal Fishing Detection,''
\textit{Int.\ J.\ Res.\ Appl.\ Sci.\ Eng.\ Technol. (IJRASET)}, vol.~13, no.~VI, 2025.

\bibitem{kim2024vessel}
J.~Kim et al.,
``Vessel Trajectory Classification using Transfer Learning with Deep Convolutional Neural Networks,''
\textit{PLOS ONE}, vol.~19, no.~3, e0298123, 2024.

\end{thebibliography}

\end{document}
